\input header
\input ras
\def\rot{{\rm rot}}
\def\Rot{{\rm Rot}}

\footline={\hss\tenrm 12-\folio\hss}\pageno=1

\centerline{\bf STRUCTURE OF THE FATOU SET}\bigskip

\centerline{\bf \S12. Herman Rings.}\medskip

The next two sections will be surveys only, with no proofs for several
major statements. This section
will describe a close relative of the Siegel disk.\smallskip

{\bf Definition.} A component \[U\] of the Fatou set \[\hat\C\ssm J(f)\] is
called a
{\bit Herman ring} if \[U\] is conformally isomorphic to some annulus
\[{\cal A}_r=\{z:1<|z|<r\}\], and if \[f\] (or some itterate of \[f\])
corresponds to an irrational rotation of this annulus.

{\bf Remark.} These objects are sometimes called
``Arnold-Herman rings'', since the existence of some examples would follow 
easily from Arnold's work in 1965.
Siegel disks and Herman rings are often collectively called
``{\bit rotation domains}''.

There are two known methods for constructing Herman rings. The original
method, due to Herman, is based on a careful analysis of real analytic
diffeomorphisms of the circle, as first studied by Arnold. An alternative
method, due to Shishikura, uses quasiconformal
surgery to cut and paste together two Siegel disks in order to fabricate
such a ring.

The original method can be outlined as follows. (Compare [S1], [He1].)
First a
number of definitions. If \[f:\R/\Z\to \R/\Z\] is an orientation
preserving homeomorphism, then we can lift to a homeomorphism \[F:\R\to\R\]
which satisfies the identity \[F(t+1)=F(t)+1\], and
is uniquely defined up to addition of an integer constant.\smallskip

{\bf Definition.} Following Poincar\'e, the {\bit rotation number}
of the lifted map \[F\] is defined to be the real number
$$	\Rot(F)\;=\;\lim_{n\to\infty}\; {F^{\circ n}(t_0)\over n}	$$
for any constant \[t_0\],
while the rotation number \[\;\rot(f)\in\R/\Z\;\] of the circle map \[f\]
is the residue class of \[\Rot(F)\] modulo 1.\smallskip

It is well known that this construction is well defined, and invariant
under orientation preserving topological conjugacy, and that it has
the following properties. (Compare Coddington and Levinson, or de Melo.)

{\QP{\bf 12.1. Lemma.} {\it The homeomorphism \[f\]
has a periodic point with period \[q\] if and only if its rotation number
is rational with denominator \[q\].}\smallskip}

{\QP{\bf 12.2. Denjoy's Theorem.} {\it If \[f\] is smooth of class \[C^2\],
and if the rotation number \[\rho=\rot(f)\]
is irrational, then \[f\] is topologically conjugate to the rotation
\[t\;\mapsto\; t+\rho\;\;(\mod 1)\].}\smallskip}

{\QP{\bf 12.3. Lemma.} {\it Consider a one-parameter family of lifted
maps of the form
$$	F_\alpha(t)\;=\;F_0(t)+\alpha\,.	$$
Then the rotation number \[\Rot(F_\alpha)\] increases
continuously and monotonically with \[\alpha\], increasing by \[+1\] as
\[\alpha\] increases by \[+1\,.\;\;(\]However, this dependence is not strictly
monotone. Rather, there is an interval of constancy corresponding to each
rational value of \[\Rot(F_\alpha)\,.\,)\]}\medskip}
\eject

In the real analytic case, Denjoy's Theorem has an analog which
can be stated as follows. Recall from \S8 that a real number \[\xi\]
is said to be {\bit Diophantine\/} if there exist a (large) number \[n\]
and a (small) number \[\epsilon\] so that the distance of \[\xi\] from
every rational number \[p/q\] satisfies \[|\xi-p/q|>
\epsilon/q^n\]. The following
was proved in a local version by Arnold, and
sharpened first by Herman and then by Yoccoz.

{\QP{\bf 12.4. Herman-Yoccoz Theorem.} \it If \[f\] is a real analytic
diffeomorphism of \[\R/\Z\] and if the rotation number \[\rho\] is
Diophantine, then \[f\] is real analytically conjugate
to the rotation \[\;t\;\mapsto\; t+\rho\;\;(\mod 1)\].\medskip}

I will not attempt to give a proof. (In the \[C^\infty\] case, Herman
and Yoccoz prove a\break corresponding
if and only if statement: {\it Every \[C^\infty\] diffeomorphism
with rotation number \[\rho\] is
\[C^\infty$-conjugate to a rotation if and only if \[\rho\] is
Diophantine.}) % However, as far as I know,
%no such sharp result is known in the real analytic case.
\smallskip

Next we will need the concept of a {\bit Blaschke product}. (Compare Problems
1-2 and 5-1.) Given any constant \[a\in\hat\C\]
with \[|a|\ne 1\], it is not difficult to show that there is one and only
one fractional linear transformation \[z\mapsto\beta_a(z)\] which maps the
unit circle \[\partial D\] onto itself fixing the base point \[z=1\],
and which maps \[a\] to \[\beta_a(a)=0\]. For example \[\beta_0(z)=z\],
\[\beta_\infty(z)=1/z\], and in general
$$	\beta_a(z)\;=\; {1-\bar a\over 1-a}\cdot{z-a\over 1-\bar a z}	$$
whenever \[a\ne\infty\].
If \[|a|<1\], then \[\beta_a\] preserves
orientation on the circle, and maps the unit disk into itself. On the other
hand, if \[|a|>1\], then \[\beta_a\] reverses orientation on \[\partial D\]
and maps \[D\] to its complement.

{\QP{\bf 12.5. Lemma.} \it A rational map of degree \[d\]
carries the unit circle into itself if and only if it
can be written as a {\bit ``Blaschke product''}
$$	f(z)\;=\;e^{2\pi it}\beta_{a_1}(z)\cdots\beta_{a_d}(z)	\eqno (*) $$
for some constants \[e^{2\pi it}\in\partial D\] and \[a_1\,,\,\ldots\,,\,a_d
\in\hat\C\ssm \partial D\].\medskip}

Here the \[a_i\] must satisfy the conditions that \[a_j\bar a_k\ne 1\]
for all \[j\] and \[k\]. For if \[a\bar b=1\], then a brief computation shows
that \[\beta_a(z)\beta_b(z)\equiv 1\]. Evidently the expression in 12.5 is
unique, since the constants \[\;e^{2\pi it}=f(1)\;\] and \[\;\{a_1\,,\,\ldots
\,,\,a_d\}\,=\,f^{-1}(0)\;\] are uniquely determined by \[f\]. The proof of
12.5 is not difficult: Given \[f\], one simply chooses any solution to the
equation \[f(a)=0\], then divides \[f(z)\] by \[\beta_a(z)\] to obtain
a rational map of lower degree, and continues inductively.\QED

Such a Blaschke product carries the unit disk into itself if
and only if all of the \[a_j\] satisfy \[|a_j|<1\]. (Compare Problems 5-1,
12-3.) However, we will rather be
interested in the mixed case, where some of the \[a_j\] are inside the unit
disk and some are outside.

{\QP\bf 12.6. Theorem. {\it For any odd degree \[d\ge 3\] we can choose a Blaschke
product \[f\] of degree \[d\] which carries the unit circle \[\partial D\] into
itself by a diffeomorphism with any desired rotation number \[\rho\]. If
this rotation number \[\rho\] is Diophantine, then \[f\] possesses a Herman
ring.}\medskip}\break

{\bf Proof Outline.} Let 
\[d=2n+1\], and choose the \[a_j\] so that \[n+1\] of them are close to
zero while
the remaining \[n\] are close to \[\infty\]. Then it is easy to check that
the Blaschke product \[z\mapsto\beta_{a_1}(z)\cdots\beta_{a_d}(z)\] is
\[C^1$-close to the identity map on the unit circle \[\partial D\].
In particular, it induces an orientation
preserving diffeomorphism of \[\partial D\]. Now multiplying by
\[e^{2\pi it}\] and using 12.3, we can adjust the rotation number to be any
desired constant. If this rotation number \[\rho\]
is Diophantine, then there
is a real analytic diffeomorphism \[h\] of \[\partial D\] which conjugates
\[f\] to the rotation \[z\mapsto e^{2\pi i\rho}z\]. Since \[h\] is real
analytic, it extends to a complex analytic diffeomorphism on some small
neighborhood of \[\partial D\], and the conclusion follows.\QED

As an example, Figure 12 shows the Julia set for the cubic rational map\break
\[f(z)\,=\,e^{2\pi it}z^2(z-4)/(1-4z)\;\] with zeros at \[0,\,0,\,4\], where
the constant
\[t=.6151732\cdots\] is adjusted so that the rotation number will be equal to
\[(\sqrt 5 -1)/2\].
There is a critical point near the center of this picture, with a
Herman ring to its left, surrounding the
superattractive basin about the origin in the left center.
This is the simplest kind
of example one can find, since Shishikura has shown that such a ring can
exist only if the degree \[d\] is at
least three, and since it is easy to check that a polynomial map cannot
have any Herman ring. (Problem 12-1 or \S17.1.)
\midinsert
\vfill
\insertRaster herm-ring.ps pixels 960 by 576 scaled 400
\bigskip\bigskip
\centerline{Figure 12. Julia set for a cubic rational map possessing
a Herman ring.}
\vfill
\endinsert

\eject

The rings constructed in this way are very special in that they are symmetric
about the unit circle, with \[f(1/\bar z)=1/\bar f(z)\]. However
Herman's original construction, based on work of Helson and Sarason,
was more flexible.
% Furthermore, the
%rotation number is necessarily Diophantine. 
Shishikura's more general
construction also avoids the need for symmetry. Furthermore
Shishikura's construction makes it clear that the possible rotation
numbers for Herman rings are exactly the same as the possible
rotation numbers for Siegel disks. In particular, any number satisfying
the Bryuno condition of \S8.4 can occur. The idea is roughly
that one starts with two rational maps having
Siegel disks with rotation numbers \[+\rho\] and \[-\rho\] respectively.
One cuts out a small concentric disk from each, and then glues the resulting
boundaries together. After making corresponding modifications at each of the
infinitely many iterated pre-images of each of the Siegel disks, Shishikura
applies
the Morrey-Ahlfors-Bers Measurable Riemann Mapping Theorem in order to
conjugate the resulting topological picture to an actual rational map.\medskip

Although Herman rings do not contain any critical
points, none-the-less they are closely associated with critical points.

{\QP\bf 12.7. Lemma. {\it If \[U\] is a
Herman ring, then every boundary point of \[U\]
belongs to the closure of the orbit of some critical point. The boundary
\[\partial U\]
has two connected components, each of which is an infinite set.}\medskip}

\noindent The proof is almost identical to the proof of 11.4.\QED % \medskip
\lln

{\bf Problem 12-1.} Using the maximum modulus principle, show that no
polynomial map can have a Herman ring.\smallskip


{\bf Problem 12-2.} For any Blaschke product \[f:\hat\C\to\hat\C\] show that
\[z\] is a critical point of \[f\] if and only if \[1/\bar z\] is a critical
point,
and show that \[z\] is a zero of \[f\] if and only if \[1/\bar z\] is a pole.
\smallskip

{\bf Problem 12-3.} A holomorphic map \[f:D\to D\] is said to be {\bit
proper} if the inverse image of any compact subset of \[D\] is compact.
Show that any proper holomorphic map from \[D\] onto itself can be expressed
uniquely as a Blaschke product $(*)$, with \[a_j\in D\].\smallskip

{\bf Problem 12-4.} Show that the rotation number \[\rot(f)\], as well as its
continued fraction expansion, can be deduced directly from the cyclic order
relations on a single orbit, without passing to the universal covering, as
follows. Let \[0=t_0\mapsto t_1\mapsto t_2\mapsto\cdots\] be the orbit of
zero, where we can choose representatives modulo 1 so that \[t_1\le t_j<
t_1+1\] for all \[j\]. Let us define \[t_j\] to be a ``closest return on the
left'' if \[t_j<0\], and if none of the numbers \[t_k\] with \[k<j\]
belong to the open interval \[(t_j\,,\,0)\]. Similarly, \[t_j\] is a ``closest
return on the right'' if the interval \[(0,t_j)\] does not contain any
\[t_k\] with \[k<j\]. If the sequence \[t_1\,,\,t_2\,,\,\ldots\] contains
first \[n_1\] closest returns on the left then \[n_2\] closest returns on
the right, and so on, show that \[\;	\rot(f)\;=\;1/(n_1+(1/n_2+1/(n_3
+\cdots)))\]. (Compare Appendix C.)\smallskip
\vfil\eject % \bigskip\bigskip

\footline={\hss\tenrm 13-\folio\hss}\pageno=1

\centerline{\bf \S13. The Sullivan Classification of Fatou Components.}
\medskip

The results in this section are due in part to Fatou and Julia, but
with very major contributions by Sullivan.\smallskip

By a {\bit Fatou component}
we will mean any connected component of the Fatou set\break \[\hat\C\ssm J(f)\].
Evidently \[f\] carries each Fatou component \[U\] onto some Fatou
component \[U'\] by a proper holomorphic map. First consider the special case
\[U=U'\].

{\QP\bf 13.1. Theorem. {\it If \[f\] maps the Fatou component \[U\] onto
itself, then there are just four possibilities, as follows. Either \[U\]
is the immediate attractive basin of an attracting fixed point, or of one
petal of a parabolic fixed point, or else \[U\] is a Siegel disk or
Herman ring.}\medskip}

Here we are lumping together the case of a superattracting fixed point, with
multiplier \[\lambda=0\], and the case of an ordinary attracting fixed
point, with \[\lambda\ne 0\]. Note that immediate attractive basins
always contain critical points by 10.2 and 10.3, while rotation domains
(that is Siegel disks and Herman rings)
evidently cannot contain critical points.\smallskip

Much of the proof of 13.1 has already been carried out in \S4. In fact,
according to 4.3 and 4.4, a priori there are just four possibilities.
Either:

(a) \[U\] contains an attractive fixed point;

(b) all orbits in \[U\] converge to a boundary fixed point;

(c) \[f\] is an automorphism of finite order; or

(d) \[f\] is conjugate to an irrational rotation of a disk, punctured
disk, or annulus.

\noindent In Case (a) we are done. Case (c) cannot occur, since our
standing hypothesis that the degree is two or more guarantees that there
are only countably many periodic points. In Case (d) we cannot have a
punctured disk, since the puncture point would have to be a fixed point
belonging to the Fatou set, so that we would just have a Siegel disk
with its center point incorrectly removed. Thus, in order to prove 13.1,
we need only show that the boundary fixed point  in Case (b)
must be parabolic. But
this boundary fixed point certainly cannot be an attracting point or a
Siegel point, since it belongs
to the Julia set. Furthermore, it cannot be repelling, since it attracts
all orbits in \[U\]. The only other possibility, which we must exclude,
is that it might be a Cremer point.\smallskip

The proof will be based on the following statement, which is
due to Douady and Sullivan. (Compare Sullivan [S1], or Douady-Hubbard [DH2,
p. 70]. For a more classical alternative, see Lyubich [L1, p. 72].)
Let
$$	f(z)\;=\;\lambda z+a_2z^2+a_3z^3+\cdots	$$
be a map
which is defined and holomorphic in some neighborhood \[U\] of the origin,
and which has
a fixed point with multiplier \[\lambda\] at \[z=0\]. By a {\bit path
converging to the origin} in \[U\] we will mean a continuous map \[p:(0,\infty)
\to U\ssm\{0\}\], where \[(0,\infty)\] is the open interval consisting
of all positive real numbers, satisfying the condition that \[p(t)\]
tends to zero as \[t\to\infty\].
\eject

{\QP{\bf 13.2. Snail Lemma.} {\it Let \[\;p:(0,\infty)\to U\ssm\{0\}\;\]
be a path which converges to the origin, and suppose that \[f\] maps \[p\]
into itself so that \[f(p(t))=p(t+1)\]. Then either the origin is
either an attracting or a parabolic fixed point. More precisely, the
multiplier \[\lambda\] satisfies either \[|\lambda|<1\] or \[\lambda=1\].}
\medskip}

\noindent Replacing \[f\] by \[f^{-1}\], we obtain the following
completely equivalent formulation, which will be useful in \S18.

{\QP{\bf 13.3. Corollary.} {\it If \[\;p:(0,\infty)\to U\ssm\{0\}\;\]
is a path which
converges to the origin, with \[f(p(t))=p(t-1)\] for \[t>1\], then either
\[|\lambda|>1\] or \[\lambda=1\].}\medskip}

{\bf Proof of 13.2.}
By hypothesis, the orbit \[p(0)\mapsto p(1)\mapsto
p(2)\mapsto\cdots\;\] in \[\;U\ssm\{0\}\;\] converges towards the origin. Thus
the origin certainly cannot be a repelling fixed point: we must have
\[|\lambda|\le 1\]. Let us assume that \[|\lambda|=1\] with
\[\lambda\ne 1\], and show that this hypothesis leads to a contradiction.

The intuitive idea can be described as follows.
As the path \[t\mapsto p(t)\] winds closer and closer to the origin,
the behavior of the map \[f\] on \[p(t)\] is more and more dominated by the
linear term \[z\mapsto\lambda z\]. Thus \[p(t+1)\approx \lambda p(t)\]
for large \[t\],
and the image must resemble a very tight spiral as shown in Figure 13.
Draw a radial segment \[E\] joining two turns of this spiral, as shown.
Then the region \[V\] bounded by \[E\] together with a segment of the
spiral will be mapped strictly into itself by \[f\]. Therefore, by the
Schwarz Lemma, the fixed point of \[f\] at the point \[0\in V\] must be
strictly attracting; which contradicts the hypothesis that \[|\lambda|=1\].
\midinsert
\hbox to \hsize{
\hfil
\hbox{\RasterBox {snail-1.ps} {640} {480} {350}}
\hskip -1in
\hbox{\RasterBox {snail-2.ps} {640} {480} {350}}
\hfil
}
\smallskip
\centerline{Figure 13. Diagram in the \[z$-plane (left) and in the
\[Z=\log(z)\] plane (right).}
\endinsert % \vfil\eject

In order to fill in the details of this argument, it is convenient
to set \[Z=\log z\], and to lift \[p\] to a continuous path
\[Z=P(t)=X(t)+iY(t)\], with \[e^{P(t)}=p(t)\]. Then evidently as \[t\to\infty\] we have
\[X(t)\to -\infty\] and
$$	P(t+1)\;=\;P(t)+ic+o(1)	\;.$$
Here \[ic\] is a pure imaginary constant with
\[e^{ic}=\lambda\], and \[o(1)\] stands for a remainder term which tends to
zero. Similarly, we can lift \[f\] to a map \[Z\mapsto F(Z)\]
which is defined on some left half-plane  \[X<{\rm constant}\], and which
satisfies
$$	F(Z)\;=\;Z+ic+o(1)	$$
as the real part \[X={\cal R}(Z)\] tends to \[-\infty\], with the same
non-zero constant \[\,ic\]. Suppose, to fix our ideas, that \[c>0\]. Then,
for \[Z=X+iY\] within some half-plane \[X<X_0\], we may assume that \[F\]
is univalent, and that the
imaginary part \[Y\] increases by at least \[c/2\] under each iteration
of \[F\]. Choose \[t_0\] large enough so that \[X(t)<X_0\] for \[t\ge t_0\].
Then the path \[t\mapsto P(t)\] cuts the half-plane \[Y>Y(t_0)\] into two
or more connected components. Let \[V_0\] be that component whose intersection
with each horizontal line is unbounded to the left. Thus \[V_0\] is
contained in the quarter-plane \[X<X_0\,,\;Y>Y(t_0)\]. Evidently \[F\] maps
\[V_0\] diffeomorphically into itself. Furthermore, the imaginary part of
\[X+iY\] increases by at least \[c/2\] under each iteration, hence the real
part must decrease towards \[-\infty\] under iteration.

Let \[V=\exp(V_0)\cup\{0\}\] be the corresponding open set in the \[z$-plane.
Then it follows that \[f\] maps \[V\] into itself,
and that every orbit of \[f\] in \[V\] converges towards the origin.
Since the open set \[V\] is Hyperbolic, it follows by the
Schwarz Lemma that the origin is an attractive fixed point; which
contradicts our hypothesis that \[|\lambda|=1\].\QED

To prove 13.3, we simply note that the orbit
$$	\cdots\mapsto p(2)\mapsto p(1)\mapsto p(0)	$$
is repelled by the origin, so the multiplier \[\lambda\] cannot be zero.
Hence \[f^{-1}\] is defined and holomorphic near the origin. Applying
13.2 to \[f^{-1}\], the conclusion follows.\QED

{\bf Proof of 13.1.} Let \[U\] be a Fatou component which is mapped into
itself by \[f\] in such a way that all orbits converge to a boundary fixed
point \[w_0\]. Choose any base point \[z_0\] in \[U\], and choose any
path \[p:[0,1]\to U\] from \[p(0)=z_0\] to \[p(1)=f(z_0)\]. Extending for
all \[t\ge 0\] by setting \[p(t+1)=f(p(t)\], we obtain a path in \[U\]
which converges to the boundary point \[w_0\] as \[t\to\infty\]. Therefore,
according to 13.2, the fixed point \[w_0\] must be either parabolic or
attracting. But \[w_0\] belongs to the Julia set, and hence cannot be
attracting.\QED\smallskip
%  \advancepageno

Thus we have classified the Fatou components which are mapped onto themselves
by \[f\]. There is a completely analogous description of Fatou components
which cycle periodically under \[f\]. These are just the Fatou components
which are fixed by some iterate of \[f\]. Hence each such component is
either

(1) the immediate attractive basin for some attracting periodic  point,

(2) the immediate attractive basin for some petal of a parabolic periodic
point,

(3) one member of a cycle of Siegel disks, or

(4) one member of a cycle of Herman rings.

\noindent In Cases (3) and (4), the topological type of the domain \[U\]
is specified by this description.
In Cases (1) and (2) it can be described as follows.

{\QP{\bf 13.4. Lemma.} \it Every immediate attractive basin is either simply
connected or infinitely connected.\medskip}

{\bf Proof.} If \[U\] were a non-simply-connected region of finite
connectivity, then it would have a finite Euler characteristic \[\chi(U)
\le 0\]. By 10.2 or 10.3, f maps \[U\] onto itself by a branched covering
map with at least
one branch point, and hence with degree \[n\ge 2\]. But the precise
number of branch points, counted with multiplicity, is equal to
\[(n-1)\chi(U)\le 0\] by the Riemann-Hurwitz formula of \S5.1. This
contradiction completes the proof.\QED	

Sullivan showed that there can be at most a finite number of such periodic
Fatou components. {\it In fact, according to Shishikura, there can be at most 2d-2
distinct cycles of Fatou components.}\smallskip

In order to complete the picture, we need the fundamental theorem
that there are no ``wandering'' Fatou components. (Compare [S2], [C].)

{\QP{\bf 13.5. Theorem of Sullivan.} \it Every Fatou component \[U\] is
eventually periodic. That is, there necessarily exist integers \[n\ge 0\]
and \[p\ge 1\] so that the \[n$-th forward image \[f^{\circ n}(U)\] is
mapped onto itself by \[f^{\circ p}\].\medskip}

Thus every Fatou component is a preimage, under some iterate of \[f\],
of one of the four types described above.
The proof, by quasiconformal deformation, is beyond the scope of these notes.
Roughly speaking, if a wandering Fatou component were to exist, then one
could construct an infinite dimensional space of deformations, all of which
would have to be rational maps of the same degree. But the space of
rational maps of fixed degree is finite dimensional.\smallskip

Recall from \S11.6
that \[f\] is {\bit post critically finite\/} (or a {\bit Thurston map\/})
if every critical orbit is finite. Combining 13.1 and
13.4 with 11.6 and 12.7, we easily obtain the following. (Compare Problem 5-7.)

{\QP{\bf 13.6. Corollary.} \it If a post critically finite
rational map has no superattractive
periodic orbit, then its Julia set is the entire sphere \[\hat\C\].\medskip}

\noindent We will give a different proof for this statement in 14.6.\lln
% \eject
We can sharpen the defining property of the Fatou set as follows.\smallskip

{\bf Problem 13-1.} If \[V\] is a connected open subset of the Fatou
set \[\hat\C\ssm J\], show that the set of all limits of
successive iterates \[f^{\circ n}|V\] as \[n\to\infty\] is either
(1) a finite set of constant maps from \[V\] into an attracting or
parabolic periodic orbit, or (2) a one-parameter family of maps,
consisting of all compositions
\[R_\theta\circ f^{\circ k}|V\] where \[f^{\circ k}\] is some fixed
iterate with values in a rotation domain and \[R_\theta\] is the rotation
of this domain through angle \[\theta\].
\bigskip




\vfil\eject
\footline={\hss\tenrm 14-\folio\hss}\pageno=1

\centerline{\bf \S14. Sub-hyperbolic and hyperbolic Maps.}\bigskip

This section will discuss two important classes of rational maps. The
exposition will be based on Douady-Hubbard [DH2].\smallskip

First some standard definitions from the theory of smooth dynamical systems.
Let \[f:M\to M\] be a \[C^1$-smooth map from a smooth
Riemannian manifold to itself, and let\break \[X\subset M\] be a compact
\[f$-invariant
compact subset, \[f(X)\subset X\]. Let\break \[Df_x:T_xM\to T_{f(x)}M\],
or briefly \[Df\], be the {\bit first
derivative map} at \[x\], that is the induced linear map on the tangent
space at
\[x\], and let \[\|v\|\] be the Riemannian norm of a vector \[v\in T_xM\].
By definition, the map \[f\] is {\bit expanding}
on \[X\] if the length of any tangent vector at a point of \[X\]
expands exponentially under iteration of \[Df\], that is, if there
exist constants \[c>0\] and \[k>1\] so that
$$	\|Df^{\circ n}(v)\|\;\ge\; ck^n\|v\|	$$
for every \[x\in X\] and \[v\in T_xM\], and every \[n\ge 0\]. Since \[X\]
is compact, a completely
equivalent requirement is that there exists some fixed \[n\ge 1\] so that
$$	\|Df^{\circ n}(v)\|\;>\;\|v\|	$$
for all non-zero tangent vectors \[v\] over \[X\]. Similarly, \[f\] is {\bit
contracting} on \[X\] if there are constants
\[c>0\] and \[k<1\] so that \[\|Df^{\circ n}(v)\|\;\le\;ck^{n}\|v\|\,\].
It is not difficult to check that
these conditions do not depend on the particular choice of Riemannian metric.
\smallskip

In the higher dimensional case, a map is called {\bit hyperbolic} on \[X\]
if the tangent space of \[M\]
restricted to \[X\] splits as the Whitney sum of two sub-bundles, each
invariant under \[Df\], so that \[Df\] is expanding on one sub-bundle and
contracting on the other. In the
one-dimensional case, this concept simplifies as follows.

{\bf Definition:} Let \[f\]
be a holomorphic map from a Riemann surface to itself, and let \[X\] be a
compact \[f$-invariant subset. The map \[f\] is
{\bit hyperbolic} on \[X\] if \[f\] is expanding on \[X\] or contracting
on \[X\], or if \[X\] is the disjoint union
of a compact \[f$-invariant subset \[X^+\] on which \[f\] is expanding
and a compact \[f$-invariant subset \[X^-\] on which \[f\] is contracting.

{\QP{\bf 14.1. Theorem.} {\it For a rational map \[f:\hat\C\to\hat\C\] of
degree \[d\ge 2\], the following two conditions are equivalent:

$(1)$ \[f\] is expanding and hence hyperbolic on its Julia set \[J\].

$(2)$ The forward orbit of each critical point of \[f\] converges
towards some attracting periodic orbit.}\medskip}
%$(4)$ The {\bit postcritical set} \[P\], that is the union of the forward
%orbits of the critical points, has closure \[\bar P\] which is disjoint
%from the Julia set \[J\].

Map satisfying these conditions are briefly called {\bit hyperbolic}
rational maps.  As examples, Figures 1a, 1b, 1d and 5 show the Julia sets of
hyperbolic maps, but the other figures do not.
{\bf Caution:} This use of the word ``hyperbolic'' has
nothing to do with the concept of ``Hyperbolic Riemann surface'',
which we write with a capital {\bit H}. (Compare \S2.)\eject % \smallskip

The proof that \[(1)\Rightarrow (2)\] proceeds as follows. We suppose
given some smooth Riemannian metric on \[\hat\C\] so that the lengths of
tangent vectors to \[\hat\C\] at points of \[J\] increase exponentially under
iteration of \[Df\]. After replacing \[f\] by some fixed iterate \[f^
{\circ p}\] (or alternatively after carefully modifying the metric), we may
assume that \[\|Df(v)\|\ge k\|v\|\] for all tangent vectors at points
within some neighborhood \[V\] of
\[J\], where \[k>1\]. In particular, there can be no critical points within
this neighborhood \[V\].

{\QP{\bf 14.2. Lemma.} \it If there exists such an expanding metric in a
neighborhood of the\break Julia set, then every orbit outside of the Julia set
converges towards an attracting periodic orbit.\medskip}

{\bf Proof.} We first show that there exists a constant
\[\epsilon>0\] so that the neighborhood \[N_\epsilon\], consisting of all
points at Riemannian distance less than \[\epsilon\] from \[J\], has the
following property:

{\QP \it If we iterate the mapping \[f\] starting at any point which
belongs to the neighborhood \[N_\epsilon\] but not to \[J\], then the
distance from \[J\] will be multiplied by
\[k\] or more at each iteration until we leave the set \[N_\epsilon\].
Furthermore, no point outside of \[N_\epsilon\] maps into \[N_\epsilon\].
\smallskip}

\noindent
In fact we need only choose \[\epsilon\] small enough so that the neighborhood
\[N_\epsilon\] is contained in \[V\], but disjoint from the compact set
\[f(\hat\C\ssm V)\]. It then follows that \[f^{-1}(N_\epsilon)\subset V\]. In
fact we will prove the stronger statement that \[f^{-1}(N_\epsilon)\subset
N_{\epsilon/k}\subset N_\epsilon\]. For if \[f(z)\in N_\epsilon\], then
\[f(z)\] can be joined
to \[J\] within \[N_\epsilon\] by a geodesic of length equal to the
Riemannian distance \[\rho(f(z),J)<\epsilon\]. We can then lift this back to
a smooth curve which joins \[z\] to \[J\] within \[V\]. The length of
this pulled back curve is evidently \[\le\rho(f(z),J)/k<\epsilon/k\].

Since \[f^{-1}(N_\epsilon)\subset N_\epsilon\], it follows that the
complementary compact set \[L=\hat\C\ssm N_\epsilon\] satisfies \[f(L)
\subset L\].
We have shown that every orbit which starts outside of \[J\] must eventually
be absorbed by this invariant set \[L\].
Since \[L\] is a compact subset of the Fatou set
\[\hat\C\ssm J\], it must be covered by finitely many of the connected components
of \[\hat\C\ssm J\]. Let \[U_1\,,\,\cdots\,,\, U_p\] be those components
of \[\hat\C\ssm J\] which intersect \[L\]. Since \[f(L)\subset
L\], it follows that \[f\] carries each one of these \[U_i\] onto some
\[U_j\]. As we follows such an orbit, we must eventually reach a component
\[U_j\] which is mapped onto itself by some iterate \[f^{\circ q}\].

Now let us examine the four possibilities for such a map \[f^{\circ q}:U_j\to
U_j\], as listed in Theorem 3.1.
This map cannot be a rotation of \[U_j\], or an automorphism
of finite order, since all orbits in \[U_j\]
are eventually absorbed by the compact subset \[L\cap U_j\]. Similarly,
orbits cannot diverge to infinity with
respect to the \Poincare metric on \[U_j\]. Thus the only
possibility is that all orbits of \[f^{\circ q}\] within \[U_j\]
converge to an attracting fixed point of \[f^{\circ q}\]. This proves
the Lemma, and hence completes the proof that
\[(1)\Rightarrow (2)\].\smallskip

{\bf Proof that \[(2)\Rightarrow (1)\].} Let \[P\] be the {\bit post-critical
set}, that is the union of the forward orbits of the critical points.
Then Condition (2) clearly implies that the closure \[\bar P\] is disjoint
from the Julia set. Conversely, if \[\bar P\cap J=\emptyset\], then we
will construct an expanding metric near the Julia set. Let
\[U\] be the open set \[\hat\C\ssm \bar P\]. Since
\[f(\bar P)\subset\bar P\], we have \[f^{-1}(U)\subset U\].
There are no critical points in \[U\], so it follows that \[f^{-1}\]
lifts to a single valued analytic function from the universal covering
surface \[\tilde U\] into itself. But \[U\] contains the Julia set, which
contains at least one repelling fixed point. Therefore this inverse map
on \[\tilde U\] can be chosen so as to contain an attracting fixed point.
If \[\tilde U\] is isomorphic to the unit disk, then this implies that
\[f^{-1}\] strictly contracts the \Poincare metric on \[\tilde U\], so that
\[f\] must strictly increase the \Poincare metric on \[U\]. This proves that
\[f\] is expanding on the compact subset \[J\subset U\].
On the other hand, if \[\tilde U\] is not isomorphic to \[D\], then the
complement \[\bar P=\hat\C\ssm U\] can contain at most two points. It is then
easy to check that \[f\]
must be conjugate to a map of the form \[z\mapsto z^{\pm d}\].
In this case, clearly \[f\] is expanding with respect to the Euclidian metric
on \[U\]. Thus \[(2)\Rightarrow (1)\], which completes the proof.\QED

One invariant set of particular interest is the ``non-wandering set''
of \[f\]. By definition, a point is {\bit wandering} if it has a
neighborhood \[U\] so that the forward images \[f(U)\,,\, f^{\circ 2}(U)\,,\,
\ldots\] are all disjoint from \[U\]. Otherwise, it is {\bit non-wandering}.
The closed set consisting of all non-wandering points is called
the {\bit non-wandering set} \[\Omega=\Omega(f)\]. The map \[f\] is said to
satisfy Smale's {\bit Axiom A} if \[f\] is
hyperbolic on \[\Omega\], and if \[\Omega\] is precisely equal to the
closure of the set of periodic points.

{\QP{\bf 14.3. Corollary.} \it A rational map satisfies Axiom A
if and only if it is hyperbolic.\medskip}

The proof is straightforward,
since we know by \S11 that the repelling periodic points are dense in the
Julia set. Details will be left to the reader. (Compare Problem 14-1.)
\QED

{\bf Remark.} These hyperbolic maps have other extremely important properties.
If \[f\] is hyperbolic, then every nearby map is also hyperbolic. Furthermore,
according to Ma\~n\'e, Sad, Sullivan, and also Lyubich, the Julia set \[J(f)\]
varies continuously under a deformation of \[f\] through hyperbolic maps.
In the non-hyperbolic case, a small change in \[f\] may well lead to a drastic
alteration of \[J(f)\]. It is generally conjectured that every rational map
can be approximated arbitrarily closely by a hyperbolic map.
\medskip


Douady and Hubbard, using ideas of Thurston, also consider a wider class
of mapping which they call {\bit sub-hyperbolic}. The only change in the
definition is that the Riemannian metric is now allowed to have a finite
number of relatively mild singularities,
like that of the metric \[|d\root n\of z |\] at the origin. Such singularities
have the property that the Riemannian distance between points still tends to
zero as the points approach each other.\smallskip

{\bf Definition.} A conformal metric on a Riemann surface, with the
expression  \[\gamma(z)|dz|\] in terms of a local uniformizing parameter
\[z\], will be called an {\bit orbifold metric} if the function \[\gamma(z)\]
is smooth and non-zero except at a locally finite collection of points
\[a_1\,,\,a_2\,,\,\ldots\;\]
where it blows up in the following special way. There should be
integers \[\nu_j\ge 2\] called the {\bit ramification indices} at the points
\[a_j\] with the following property. If we take a local branched covering
by setting \[z(w)=a_j+w^{\nu_j}\], then the induced metric \[\gamma(z(w))
|(dz/dw)dw|\] on the \[w$-plane should be smooth and non-singular throughout
some neighborhood of the origin.

The rational map \[f\] is {\bit
sub-hyperbolic} if it is expanding with respect to some orbifold metric
on a neighborhood of its Julia set.
\eject

{\QP{\bf 14.4. Theorem.} \it A rational map is sub-hyperbolic if and only if:

$(\alpha)$ every critical orbit in the Julia set is eventually periodic,
and

$(\beta)$ every critical orbit outside of the Julia set converges to an
attracting periodic orbit.\medskip}

(Compare 11.7.) As examples, Figures 1-5 show sub-hyperbolic maps,
but 8-10, 12, 17 do not.

The proof in one direction is completely analogous to
the proof above. If \[f\] is expanding with respect to some orbifold metric,
defined near the Julia set, then, just as in 14.2 above, every orbit
outside the Julia
set \[J\] will be pushed away from \[J\], and can only converge to an
attractive periodic orbit. On the other hand, if \[\omega\] is a critical
point inside the Julia set, then every post-critical
point \[f^{\circ n}(\omega)\,,\; n>0\], must be one of the singular
points  \[a_j\] for our
orbifold metric, since the map \[f^{\circ n}\] is critical, and yet is required
to be expanding, at the point \[\omega\]. Thus the forward
orbit of \[\omega\] cannot have any limit points, and hence must be finite.
\smallskip

For the proof in the other direction, we must introduce the concept of the
``universal branched covering'' of an orbifold. Compare Appendix E.
For our purposes, an
{\bit orbifold} \[(S,\,\nu)\] will just mean a Riemann surface
\[S\], together with a locally finite collection of marked points \[a_j\],
each of which is assigned
a {\bit ramification index} \[\nu_j=\nu(a_j)\ge 2\] as above. For any
point \[z\] which is not one of the \[a_j\] we set \[\nu(z)=1\].\smallskip

{\bf Definition.} To every rational
map \[f\] which satisfies Conditions ($\alpha$) and ($\beta$) of Theorem
14.4, we assign
the {\bit canonical orbifold} \[(S,\,\nu)\] as follows. As underlying
Riemann surface \[S\] we take the Riemann sphere \[\hat\C\] with all
attracting periodic orbits removed. As ramified points \[a_j\] we take
all strictly {\bit post-critical} points, that is, all points which have
the form \[a_j=f^{\circ n}(\omega)\] for some \[n>0\], where \[\omega\] is
a critical point of \[f\]. Using ($\alpha$) and ($\beta$), we see easily that this
collection of points \[a_j\] is locally finite in \[S\] (although perhaps
not in \[\hat\C\]). In order to specify the corresponding ramification indices
\[\nu_j=\nu(a_j)\], we will need another definition.
If \[f(z_1)=z_2\], with local power series development
$$	f(z)=z_2+c(z-z_1)^n+({\rm higher\;terms})\,,	 $$
where \[c\ne 0\] and \[n\ge 1\],
then the integer \[n=n(f\,,\,z_1)\] is called the {\bit local degree} or the
{\bit branch index} of \[f\] at \[z_1\].
{\it Now choose the \[\nu(a_j)\ge 2\] to be the smallest integers which
satisfy the following:}

{\QP{\bf Condition ($\ast$).} For any \[z\in S\],
the ramification index \[\nu(f(z))\] at the image point must be a multiple
of the product \[n(f\,,\,z)\nu(z)\]. \medskip}

\noindent
In fact, to construct these integers \[\nu(a_j)\]
we simply consider all critical pre-images
$$	f^{\circ m}(\omega)=a_j\,,\quad f'(\omega)=0\,,	$$
and let \[\nu(a_j)\] be the least common multiple of the corresponding
branch indices\break \[n(f^{\circ m}\,,\,\omega)\]. Since we have removed all
attracting periodic orbits, it is not difficult to
check that this least common multiple is finite,
and that it provides a minimal solution to the required condition $(\ast)\].
\smallskip

As in Appendix E, we consider the {\bit universal covering surface}
$$	\tilde S_\nu\;\to\; (S\,,\,\nu)	$$
for this canonical orbifold, that is the unique regular branched
covering of \[S\] which has the given \[\nu\] as ramification function,
and which is simply connected. In order
to determine the geometry of \[\tilde S_\nu\], we introduce
the {\bit orbifold Euler characteristic}
$$	\chi(S\,,\,\nu)\;=\;\chi(S)+\sum({1\over\nu(a_j)}-1)\;=\;2-k+
	\sum({1\over\nu(a_j)}-1)\,,	$$
to be summed over all ramified points, where \[k\] is the number of
attracting periodic points. According to Lemma E.1 in the Appendix,
such a universal covering nearly always exists; and the few pathological
cases where it does not exist necessarily have Euler characteristic \[\chi
(S\,,\,\nu)>0\].
Furthermore, by E.4, the surface \[\tilde S_\nu\] is either Spherical,
Hyperbolic, or Euclidean according as \[\chi(S\,,\,\nu)\] is positive,
negative, or zero.\smallskip

Still assuming Conditions ($\alpha$) and ($\beta$) of 14.4, we will
prove the following.

{\QP{\bf 14.5. Lemma.} \it This canonical orbifold has Euler
characteristic
\[\chi(S\,,\,\nu)\le 0\].\break Hence its universal covering \[\tilde
S_\nu\]
has either a unique \Poincare ~metric or a\break
Euclidean metric which is unique
up to a scale change. In either case, there is an induced orbifold metric
on \[S\], and \[f\] is expanding with respect to this orbifold metric
throughout any compact subset of \[S\].\medskip}

Evidently this will complete the proof of Theorem 14.4.
In the Euclidean case, the proof yields much more. Let \[z\] be a uniformizing
parameter on the covering surface \[\tilde S_\nu\cong\C\].

{\QP{\bf 14.6. Theorem.} \it If \[\chi(S\,,\,\nu)=0\], then \[f\] induces
a linear isomorphism\break \[z\mapsto az+b\] from the Euclidean covering space
\[\tilde S_\nu\cong \C\] onto itself. In this case, the Julia set is
either a circle or line segment, or the entire Riemann sphere. The
coefficient of expansion \[|a|\] satisfies \[|a|=d\] in the first two
cases, and \[|a|=\sqrt d\] in the last case, where \[d\] is the degree
of \[f\].\medskip}

(Caution: The coefficient \[a\] itself is not uniquely determined; for
the lifting of \[f\] to the covering surface is determined only up to
composition with a deck transformation. The deck transformations
may well have fixed points, since we are dealing with a branched covering,
but necessarily
have the form \[z\mapsto \omega z+c\] where \[\omega\] is some root of unity.)
\smallskip

{\bf Proof of 14.5.} First note that Condition ($\ast$) above is exactly
what is needed in order to assert that \[f^{-1}\] lifts, locally at least,
to a holomorphic map from the universal covering \[\tilde S_\nu\] to
itself. But since \[\tilde S_\nu\] is simply connected, there is then no
obstruction to constructing a globally lifted inverse
$$	\tilde f^{-1}\,:\, \tilde S_\nu\;\to\;\tilde S_\nu\,.	$$
The proof now divides into three cases.\smallskip

{\bf Spherical Case.} Note that the composition
$$	\tilde S_\nu\;\buildrel \tilde f^{-1}\over\longrightarrow\;
	\tilde S_\nu\;\buildrel p\over\longrightarrow\;S\;\buildrel
	 f\over\longrightarrow\;S $$
always coincides with the projection map \[p:\tilde S_\nu\to S\].
If \[\tilde S_\nu\] were Spherical, then this composition
would have a well defined
degree, which would have to be at least \[d\] times the degree of \[p\]. This
is impossible, hence the Spherical case cannot occur.\smallskip

{\bf Hyperbolic Case.} If \[\tilde S_\nu\] is Hyperbolic, then it has a
uniquely defined Poincar\'e metric, and the proof proceeds just as in the
proof that \[(2)\Rightarrow (1)\] in 14.1 above.\smallskip

{\bf Euclidean Case.} Suppose that \[\tilde S_\nu\] is Euclidean, or
equivalently that \[\chi(S\,,\,\nu)=0\].
Construct a new orbifold \[(S'\,,\,\mu)\] as follows. Let \[S'\subset S\]
be the open set \[S\] with all immediate pre-images of attracting periodic
points removed, and define \[\mu=f^*(\nu)\] by the formula
$$	\mu(z)\;=\; \nu(f(z))/n(f\,,\,z)	$$
where \[n(f\,,\,z)\] is the branch index.
Note that \[\mu(z)\ge\nu(z)\] for all \[z\in S'\]. Evidently it follows
that
$$	\chi(S'\,,\,\mu)\;\le\; \chi(S\,,\,\nu)	$$
with equality only if \[S'=S\] and \[\mu=\nu\]. But by Lemma E.2, since the
map\break \[f:(S'\,,\,\mu)\to(S\,,\,\nu)\] is a ``$d$-fold covering'' of
orbifolds, we conclude that \[f\] induces an isomorphism \[\tilde S'_\mu
\to\tilde S_\nu\] of universal
covering surfaces, and also that the Riemann-Hurwitz formula takes the form
$$	\chi(S'\,,\,\mu)\;=\;\chi(S\,,\,\nu)d\,.	$$
Thus \[\chi(S'\,,\,\mu)\] is also zero. We conclude that
\[S'=S\] and \[\mu=\nu\], so that \[f\] must lift to a (necessarily linear)
isomorphism from the Euclidean covering space \[\tilde S_\nu\] to itself.
Further details will be left to the reader. This completes the proof
of 14.4 through 14.6.\QED

As a corollary, we obtain another proof of 13.6.
We continue to assume that every
critical orbit either converges to an attracting orbit or is eventually
periodic, according as it belongs to the Fatou set or the Julia set.

{\QP{\bf 14.7. Corollary.} \it If \[f\] is sub-hyperbolic with
no attracting periodic orbits,
so that \[S\] is the entire Riemann sphere, then \[f\] is expanding on the
entire sphere, and it follows that \[J(f)\] is the entire sphere.
\medskip}

\lln

{\bf Problem 14-1.} Using the results of \S13, show that the non-wandering
set for a rational map \[f\] is the (disjoint) union of its Julia set,
its rotation domains (if any), and
its set of attracting periodic points.\smallskip

{\bf Problem 14-2.} Show that the Julia set for the rational map \[z\mapsto
(1-2/z)^{2n}\] is the entire Riemann sphere. Show that the orbifold
metric for this example is Euclidean when \[n=1\], but is Hyperbolic
for \[n>1\].\smallskip

{\bf Problem 14-3.} For any sub-hyperbolic map whose
canonical orbifold metric is\break Euclidean,
show that every periodic orbit outside of the finite post-critical set has
multiplier \[\lambda\] satisfying \[|\lambda|=d^{p/\delta}\] where
\[d\] is the degree,
\[p\] is the period, and \[\delta\] is the dimension (1 or 2) of the Julia
set.\smallskip

{\bf Problem 14-4.} A rational map \[f\] is said to be {\bit expansive}
in a neighborhood of its Julia set if there exists \[\epsilon>0\] so that,
for any two points \[x\ne y\] whose orbits remain in the neighborhood forever,
there exists some \[n\ge 0\] so that \[f^{\circ n}(x)\] and \[f^{\circ n}(y)\]
have distance greater than \[\epsilon\]. Using Sullivan's results from \S13,
show that this condition is satisfied if and only if \[f\] is hyperbolic.
(However, a map with a parabolic fixed point may be expansive on the Julia set
itself.)

\vfill
%\rightline{jwm, version of 9-2-91}
\end









